\chapter{Fonctions exponentielles et logarithmiques}
\label{workbook:chapter:06}
\section{Variation additive et variation multiplicative}
\label{workbook:section:06.01}
\begin{exerciseblock}{Variation additive et variation multiplicative}
\item Comparer les suites $u_n=100+8n$ et $v_n=100(1{,}08)^n$ après $1$, $5$ et $10$ périodes.
\item À partir d'un tableau, expliquer comment reconnaître une variation additive constante et une variation multiplicative constante.
\item Une quantité diminue de $12\%$ par année. Écrire le facteur multiplicatif et le modèle après $n$ années à partir de $Q_0$.
\end{exerciseblock}
\section{Le logarithme comme fonction réciproque}
\label{workbook:section:06.02}
\begin{exerciseblock}{Le logarithme comme fonction réciproque}
\item Convertir entre formes exponentielle et logarithmique : $2^5=32$, $10^{-3}=0{,}001$, $\log_3 81=4$.
\item Résoudre $5^x=17$ en expression exacte puis approchée.
\item Expliquer graphiquement pourquoi $y=b^x$ et $y=\log_bx$ sont symétriques par rapport à $y=x$.
\end{exerciseblock}
\section{Lois des logarithmes}
\label{workbook:section:06.03}
\begin{exerciseblock}{Lois des logarithmes}
\item Développer $\log_b\left(\dfrac{x^3\sqrt y}{z^2}\right)$ sous les hypothèses appropriées.
\item Condensez $2\ln x-\frac12\ln y+\ln3$ en un seul logarithme.
\item Montrer par un contre-exemple que $\ln(x+y)\neq\ln x+\ln y$.
\end{exerciseblock}
\section{Équations exponentielles et logarithmiques}
\label{workbook:section:06.04}
\begin{exerciseblock}{Équations exponentielles et logarithmiques}
\item Résoudre $3^{2x-1}=7$.
\item Résoudre $\ln(x-1)+\ln(x+2)=\ln 12$ en contrôlant le domaine.
\item Résoudre $2e^{0{,}4t}=15$ et interpréter si $t$ représente un temps.
\end{exerciseblock}
\section{Applications}
\label{workbook:section:06.05}
\begin{exerciseblock}{Applications}
\item Un capital de $5000\,$\$ croît de $4{,}2\%$ annuellement. Calculer sa valeur après $8$ ans et le temps de doublement.
\item Une substance a une demi-vie de $6$ h. Écrire le modèle et calculer la fraction restante après $15$ h.
\item Le pH vaut $-\log_{10}[H^+]$. Comparer les concentrations associées aux pH $3$ et $5$.
\end{exerciseblock}
\section{Interprétation d'un facteur multiplicatif}
\label{workbook:section:06.06}
\begin{exerciseblock}{Interprétation d'un facteur multiplicatif}
\item Pour $f(t)=240(1{,}03)^t$, interpréter tous les paramètres et calculer le taux de variation relatif sur une période.
\item Écrire un modèle équivalent sous la forme $Ae^{kt}$ et déterminer $k$ pour un facteur annuel $1{,}03$.
\item Une quantité est multipliée par $0{,}85$ par période. Distinguer diminution en pourcentage et perte absolue.
\end{exerciseblock}
\section{Trois façons de voir une exponentielle}
\label{workbook:section:06.07}
\begin{exerciseblock}{Trois façons de voir une exponentielle}
\item Pour $f(t)=80\cdot2^{t/5}$, décrire la formule, un tableau à pas $5$ et l'allure du graphe.
\item Prendre le logarithme de $y=Ab^x$ et montrer que $\ln y$ est affine en $x$.
\item Des données donnent $(x,y)=(0,3),(1,6),(2,12),(3,24)$. Vérifier le modèle par rapports et par logarithmes.
\end{exerciseblock}
\section{Modéliser une variation multiplicative}
\label{workbook:section:06.08}
\begin{exerciseblock}{Modéliser une variation multiplicative}
\item Une population passe de $1200$ à $1800$ en $4$ ans. Construire un modèle $P(t)=1200b^t$ et calculer $b$.
\item Un médicament décroît de $80$ mg à $30$ mg en $6$ h selon $M(t)=80e^{-kt}$. Déterminer $k$.
\item Un refroidissement suit $T(t)=T_a+(T_0-T_a)e^{-kt}$. Vérifier que $T(t)\to T_a$ et interpréter chaque terme.
\end{exerciseblock}
\section*{Synthèse du chapitre}
\begin{synthesisbox}
\item Comparer sur $20$ ans un placement simple $C_s(t)=1000(1+0{,}05t)$ et composé $C_c(t)=1000(1{,}05)^t$. Trouver quand le composé dépasse le simple de $200\,$\$.
\item Une culture triple toutes les $7$ h. Partant de $400$ cellules, quand atteint-elle $10^6$ cellules?
\item Construire à partir de deux points un modèle exponentiel, puis vérifier graphiquement que son logarithme est affine.
\end{synthesisbox}
